\input{../../lean/.lake/build/blueprint/module/Cliquefree/Blueprint}

\tableofcontents

\chapter{Conventions}

This blueprint presents the weighted independent-set bound in
Theorem~\ref{thm:existence} and the main steps of its checked Lean proof.
Its coefficient is
$c_r=(1+1152B_0r)^{-1}=\Theta(r^{-1})$.

Scalar analysis supplies a potential and its tangent-loss bound. A triangle
dichotomy and sparse neighborhood restriction rule out triangle witnesses at
a positive maximizer. The resulting triangle bound makes a deletion reward
nonnegative. Compactness and induction then establish the weighted bound.

Graphs are finite and simple, with vertex set $V$. A graph is $K_r$-free if
it has no set of $r$ pairwise adjacent vertices; a set is independent if
no two of its vertices are adjacent. The weights $s>0$ prescribe degrees,
$u\ge0$ is the objective, and $w$ denotes auxiliary weights. For any weight
vector $h$ and vertex set $A$, write $h(A)=\sum_{v\in A}h_v$. Thus
$\alpha_u(G)=\max\{u(I):I\text{ independent}\}$.
All vector inequalities are coordinatewise. The symbols $N(v)$ and
$N[v]=N(v)\cup\{v\}$ denote open and closed neighborhoods, and
$N[A]=\bigcup_{v\in A}N[v]$. Write $G[A]$ for the subgraph induced by $A$,
and $G-A=G[V\setminus A]$; weights on these subgraphs are restricted to
their vertex sets. The cross mass
$\Xi_w(A,B)=\sum_{x\in A}\sum_{y\in B\cap N(x)}w_xw_y$ counts directed
incidences; edge and triangle masses count each unordered edge or triangle once.
All logarithms are natural unless a base is shown. The absolute constant
$B_0=2^X$ is kept symbolic.

\chapter{The existence theorem}

\inputleannode{thm:existence}
\inputleannode{bp:triangle-constants}
\inputleannode{bp:constants}
\inputleannode{bp:constant-bounds}

\chapter{Weighted graphs and potentials}

\inputleannode{lean:shearer-f}
\inputleannode{bp:weighted-degree}
\inputleannode{bp:objective}
\inputleannode{lean:independence-attained}
\inputleannode{bp:extremal-benchmark}
\inputleannode{bp:graph-masses}
\inputleannode{bp:scalar-potential}
\inputleannode{bp:graph-potential}
\inputleannode{bp:residual-reward}
\inputleannode{bp:triangle-scales}
\inputleannode{bp:witness}

\chapter{Scalar analysis and lifting}

\inputleannode{lean:f-continuous}
\inputleannode{lean:f-limit}
\inputleannode{lean:f-integral}
\inputleannode{lean:optimizer-equation}
\inputleannode{lean:envelope-derivative}
\inputleannode{lean:potential-continuous}
\inputleannode{lean:xi-positive}
\inputleannode{bp:scalar-ode}
\inputleannode{bp:scalar-loss}
\inputleannode{bp:potential-comparison}
\inputleannode{lem:lift}
\inputleannode{cor:transfer}
\inputleannode{lean:maximizer-stationary}

\chapter{The triangle dichotomy}

\inputleannode{lean:pruning}
\inputleannode{lean:grand-coupling}
\inputleannode{bp:triangle-reduction}
\inputleannode{lem:triangle-test}

\chapter{Sparse restriction and stationarity}

\inputleannode{bp:disjointification}
\inputleannode{lem:sparse-restriction}
\inputleannode{lem:improvement}
\inputleannode{lean:stationary-budget}
\inputleannode{cor:stationary}

\chapter{Deletion and compactness induction}

\inputleannode{lem:deletion}
\inputleannode{lean:deletion-slack}
\inputleannode{lean:deletion-pivot}
\inputleannode{bp:maximizing-sequence}
\inputleannode{bp:boundary-decomposition}
\inputleannode{lean:compactness-reduction}
\inputleannode{lean:weighted-bound}
